% Phụ lục B: scene ranker, bộ xếp lại thứ hai trên VG ∩ COCO, phát triển song song ở repo cá nhân của một thành viên.
% Không tính vào giới hạn trang. Nguồn số: report/scene-ranker.md, từ các file chép nguyên trong
% experiments/improve/scene-ranker/personal-repo/ (commit cdb4b5b của repo cá nhân): test_eval_0070_metrics.json,
% test_eval_0069_val_metrics.json, rank_train_{0056,0057,0063,0065}_metrics.json, rank_ablation_{0058,0062,0064}_metrics.json,
% comparison_test.json.
\section{Scene ranker: khớp scene graph bằng Sinkhorn}\label{app:scene-ranker}

Phụ lục này mô tả dòng scene ranker của Bảng~\ref{tab:l1-test}. Mô hình được phát triển, train và chấm test ở một repo cá
nhân, trên cùng split Karpathy của VG $\cap$ COCO, cùng pool và cùng CLIP ViT-B/32 với repo nhóm; CLIP thuần tính trong repo
đó cho R@1 test 39,32 và MRR@50 51,99, khớp với 39,32 và 52,00 của repo nhóm. Hai repo xử lý dữ liệu độc lập nên bộ caption
lệch 3 câu ở val (10\,630) và 6 câu ở test (10\,690). Các file kết quả được chép nguyên vào
\file{experiments/improve/scene-ranker/personal-repo/}; mô tả đầy đủ ở \file{report/scene-ranker.md}.

\subsection{Mô hình}
\textbf{Hai scene graph.} Câu được FACTUAL-T5 base tách thành thực thể, thuộc tính và bộ ba, tối đa 8 thực thể và 6 quan hệ.
Ảnh dùng chú thích VG: 64 vật thể có khung lớn nhất, tối đa 4 thuộc tính mỗi vật thể, 64 quan hệ.

\textbf{Mã hoá node.} Mỗi node gộp vector CLIP text của tên, của thuộc tính, đặc trưng số (IDF; tỉ lệ khung và số lượng ở
phía ảnh) và vector concept sau GAT:
$\mathbf h^{(0)}_n = W \mathbf f_n + \overline{W_a \mathbf f_a} + W_x \mathbf x_n + \overline{\operatorname{GAT}(\mathbf z)_c}$.
GAT 2 lớp, 4 head, chạy trên một concept graph gồm 1\,687 synset xuất hiện trong ít nhất 20 ảnh train và 13\,807 cạnh (1\,486
cạnh hypernym của WordNet, 12\,114 cạnh đồng xuất hiện trên ảnh train, 207 cạnh ConceptNet 5.7). Sau đó 2 lớp attention dùng
chung cho câu và ảnh truyền tin dọc cạnh quan hệ của scene graph. Vector để khớp là
$\mathbf p_n = \operatorname{normalize}(\mathbf f_n + U \mathbf h_n)$ với $U$ khởi tạo 0.

\textbf{Khớp Sinkhorn có thùng rác.} Hai kênh: thực thể câu với vật thể ảnh, quan hệ câu với quan hệ ảnh. Mỗi kênh đưa ma
trận cosine qua 20 vòng Sinkhorn với nhiệt độ $\tau$ và điểm thùng rác $z$ học được:
$P = \operatorname{Sinkhorn}(\mathit{sim}/\tau, z)$, $S_{\text{ch}} = \sum_i w_i \sum_j P_{ij}\,\mathit{sim}_{ij}$. Mỗi phần của
câu gửi đi một đơn vị khối lượng và mỗi phần của ảnh nhận tối đa một đơn vị; phần không có đối tác đủ tốt đi vào thùng rác và
không cộng điểm, nên một vật thể của ảnh không bị tính nhiều lần.

\textbf{Điểm và huấn luyện.} Điểm đồ thị cộng hai kênh với sáu term cố định (điểm đường đi trên concept graph, thuộc tính khớp
theo từ và theo CLIP text, số lượng, độ nổi bật theo diện tích, độ giống giữa cụm danh từ và nhãn vật thể); điểm cuối là
$0{,}65\cdot\text{CLIP} + 0{,}35\cdot\text{đồ thị}$. Loss là cross-entropy của điểm cuối cộng 0,3 lần cross-entropy của riêng
điểm đồ thị. Mô hình có 709\,469 tham số; AdamW, lr 0,001, batch 64 caption, tối đa 10 epoch, chọn epoch theo MRR val.

\textbf{Chống rò rỉ.} Từ vựng concept, IDF và cạnh đồng xuất hiện chỉ đếm trên ảnh train. $\alpha$, epoch và mọi siêu tham số
chọn trên val. Test mở đúng một lần (run 0070), sau khi cùng code đánh giá chạy lại toàn bộ chuỗi trên val và khớp từng bước
với các run đã lưu (run 0069).

\subsection{Kết quả và ablation}

\begin{table}[!htb]
\centering
\caption{Scene ranker trên val và test. B1, B2: hai mốc không học (cộng IDF các concept khớp; mở rộng bằng đường đi
$\le 2$ cạnh). MRR dạng tỉ lệ, tính trên toàn gallery như repo cá nhân.}
\label{tab:app-sr-results}
\footnotesize
\setlength{\tabcolsep}{2.5pt}
\begin{tabular}{@{}llcccc@{}}
\toprule
 & Mô hình & R@1 & R@5 & R@10 & MRR \\
\midrule
\multirow{4}{*}{Val} & CLIP & 39,14 & 67,89 & 78,98 & 0,5239 \\
 & B1 & 40,20 & 68,76 & 79,99 & 0,5350 \\
 & B2 & 40,25 & 68,72 & 79,93 & 0,5352 \\
 & Scene ranker & \pmstd{49,69}{0,05} & \pmstd{76,00}{0,18} & \pmstd{85,11}{0,12} & \pmstd{0,6165}{0,0003} \\
\midrule
\multirow{4}{*}{Test} & CLIP & 39,32 & 66,38 & 77,32 & 0,5203 \\
 & B1 & 40,37 & 67,46 & 78,39 & 0,5307 \\
 & B2 & 40,41 & 67,46 & 78,40 & 0,5314 \\
 & Scene ranker & \pmstd{48,57}{0,36} & \pmstd{75,28}{0,03} & \pmstd{84,54}{0,11} & \pmstd{0,6080}{0,0021} \\
\bottomrule
\end{tabular}
\tabsrc{\file{test_eval_0070_metrics.json}, \file{rank_train_*_metrics.json}.}
\end{table}

\begin{table}[!htb]
\centering
\caption{Ablation lúc chấm của scene ranker trên val (checkpoint đã train, phá đồ thị lúc chấm), trung bình 3 seed.}
\label{tab:app-sr-ablation}
\footnotesize
\begin{tabular}{@{}lcc@{}}
\toprule
Biến thể & MRR & R@1 \\
\midrule
Mô hình đầy đủ & \pmstd{0,6165}{0,0003} & 49,69 \\
Xáo cạnh concept graph & \pmstd{0,6153}{0,0004} & 49,56 \\
Bỏ cạnh concept graph & \pmstd{0,6124}{0,0009} & 49,14 \\
Bỏ 6 term cố định & \pmstd{0,6150}{0,0008} & 49,58 \\
Bỏ kênh quan hệ & \pmstd{0,6105}{0,0013} & 48,99 \\
Xáo cạnh scene graph (câu và ảnh) & \pmstd{0,5986}{0,0015} & 47,62 \\
Bỏ cạnh scene graph (câu và ảnh) & \pmstd{0,5375}{0,0019} & 40,86 \\
Bỏ phần học được ($U\mathbf h = 0$) & \pmstd{0,5528}{0,0007} & 42,57 \\
Scene graph của ảnh khác & \pmstd{0,3892}{0,0043} & 24,96 \\
CLIP & 0,5239 & 39,14 \\
\bottomrule
\end{tabular}
\tabsrc{\file{rank_ablation_{0058,0062,0064}_metrics.json}.}
\end{table}

Trên test, từng seed hơn CLIP với khoảng tin cậy 95\% của chênh MRR (bootstrap ghép cặp) là [+0,079; +0,097],
[+0,077; +0,094] và [+0,081; +0,098] (Bảng~\ref{tab:app-sr-results}). Chấm bằng scene graph của ảnh khác làm MRR rơi dưới
CLIP, nên điểm đồ thị đọc nội dung của đúng ảnh (Bảng~\ref{tab:app-sr-ablation}). Phá cạnh lúc chấm đưa đầu vào ra ngoài phân
phối train, nên mức giảm là cận trên của vai trò cạnh. Train lại từ đầu không có cạnh nào (run 0065, một seed) đạt MRR val
0,6071, chỉ thua mô hình có cạnh 0,0097 (KTC 95\% [$-0{,}0141$; $-0{,}0055$]): cạnh đóng góp khoảng 10\% phần tăng so với
CLIP, phần còn lại đến từ việc khớp từng phần của câu với từng vật thể của ảnh. Kết luận này cùng chiều với ablation bỏ cạnh
của mô hình chính (Bảng~\ref{tab:l1-ablation-main}).

\subsection{Giải thích và 10 truy vấn}
Điểm cuối là tổng các số hạng, nên lời giải thích của một cặp (câu, ảnh) là bảng các số hạng: phần CLIP, từng thực thể và
quan hệ của câu (phần ảnh nhận nhiều khối lượng nhất, phần vào thùng rác, đóng góp vào điểm) và từng term cố định. Trên val,
seed 0: tổng các phần khớp điểm cuối với sai số tối đa $1{,}4\cdot10^{-5}$ trên 10 ứng viên đầu của mọi caption; với 10
caption có ảnh đúng ở hạng 1, xoá vật thể đóng góp nhiều nhất làm điểm giảm mạnh hơn xoá ngẫu nhiên ở 10/10 ca (trung bình
$-1{,}45$ so với $+0{,}006$), và 5/10 ca ảnh đúng mất hạng 1. Bảng~\ref{tab:app-sr-cases} là 5 truy vấn thành công (CLIP xếp
ảnh đúng ở hạng 39--45, mô hình đưa lên hạng 1) và 5 truy vấn thất bại (CLIP hạng 1, mô hình đẩy xuống hạng 9--10).

\begin{table}[!htb]
\centering
\caption{Mười truy vấn val của scene ranker. Hạng ảnh đúng theo CLIP $\to$ theo mô hình.}
\label{tab:app-sr-cases}
\scriptsize
\setlength{\tabcolsep}{2.5pt}
\begin{tabular}{@{}>{\raggedright\arraybackslash}p{3.3cm}c>{\raggedright\arraybackslash}p{3.9cm}@{}}
\toprule
Truy vấn & Hạng & Cặp khớp quyết định / nguyên nhân \\
\midrule
a device to use and work on bike fenders & 45 $\to$ 1 & bike fenders $\to$ fender; device $\to$ fender blender \\
The organizer is sitting next to the computer screen. & 45 $\to$ 1 & organizer $\to$ organizer; screen $\to$ computer \\
A large city bus is parked on the side of a street. & 42 $\to$ 1 & city bus $\to$ bus; street $\to$ street \\
A woman with a red scarf and black leather jacket walking in the city. & 40 $\to$ 1 & scarf $\to$ scarf; woman $\to$ lady \\
A mattress is leaned up against a wall & 39 $\to$ 1 & mattress $\to$ mattress; bộ ba ``mattress against wall'' \\
\midrule
A man eats a hotdog in a stadium. & 1 $\to$ 10 & Ảnh đúng thiếu ``man'', ``hot dog''; mọi thực thể vào thùng rác \\
THERE IS A WALL CLOCK ABOVE THE MANTEL PIECE & 1 $\to$ 10 & Parser tách sai ``mantel piece'' thành ``building'' \\
A motorcycle frame that has been made in to a tricycle. & 1 $\to$ 10 & Xe trong ảnh đúng gán nhãn ``cart''; câu tả một biến đổi \\
A man in the kitchen holding a beverage. & 1 $\to$ 10 & Chú thích ảnh đúng chỉ có tường và tranh \\
A zebra grazes in a cage at a fair. & 1 $\to$ 9 & Ảnh đúng thiếu nhãn ``zebra'' \\
\bottomrule
\end{tabular}
\tabsrc{\file{report/scene-ranker.md}, mục 5.}
\end{table}

Nguyên nhân chính của thất bại là chú thích thiếu trong ảnh đúng, giống mô hình chính (Mục~\ref{sec:l1-cases}); lời giải thích
chỉ ra phần nào của câu rơi vào thùng rác, nên lỗi chẩn đoán được trực tiếp. Hạn chế: lần train lại không cạnh chỉ có một
seed, và bốn thành phần của mô hình (Sinkhorn, concept graph, term cố định, loss phụ) đã được thêm vào code của nhóm dưới dạng
cờ nhưng chưa được train trên mô hình chính.
